\section{Introduction}
Routine clinical laboratory and examination data carry most of the information used to recognize chronic disease, yet they differ from the tables on which general-purpose learning methods are developed. A serum value acquires its meaning from its unit, specimen, assay and collection time; whether a test was ordered at all reflects the clinician's suspicion and is therefore informative; and each disease is studied in its own cohort, with its own feature set and label definition. A system that serves several diseases must represent this heterogeneity explicitly, predict each endpoint as accurately as the available biomarkers allow, and help decide which additional test is worth its cost.

Deep tabular models have only recently become competitive with gradient-boosted trees on such data \cite{grinsztajn2022,gorishniy2021}, and the strongest small-data learner to date is a prior-fitted transformer that performs Bayesian inference in context \cite{tabpfn,tabpfn2}. Clinical graph models inject structural priors into attention \cite{gct}, and missing-data theory shows that the Bayes-optimal predictor depends on the missingness pattern whenever values are missing not at random \cite{littlerubin,neumiss}. These ingredients have so far been studied separately.

We present PredictX-HARBOR, a multi-disease prediction system built on a typed, availability-aware biomarker representation shared by three components: a graph encoder, a heterogeneous ensemble and a test-acquisition controller. The contributions are as follows.
\begin{enumerate}\itemsep1pt
\item \textbf{HARBOR}, a reliability-conditioned biomarker graph encoder in which a value-message graph over observed biomarkers and a missingness-state graph over all biomarkers are mixed by an availability-dependent gate (Section~\ref{sec:harbor}). Unobserved biomarkers retain their concept identity, graph sparsity is regularized by attention entropy, and we prove that the $\ell_1$ attention penalty has identically zero gradient under softmax normalization (Proposition~\ref{prop:l1}).
\item \textbf{The PredictX ensemble}, an equal-weight combination of TabPFN, CatBoost, penalized logistic regression and HARBOR with a proper-score guarantee (Proposition~\ref{prop:ens}), together with cross-site training on cohorts that record the same concepts with different test availability (Section~\ref{sec:ensemble}).
\item \textbf{A value-of-information test-acquisition controller} with a non-parametric conditional surrogate and a pattern-complete predictor, applied to sequential ordering of exercise electrocardiography, fluoroscopy and thallium scintigraphy (Section~\ref{sec:voi}).
\item \textbf{An evaluation on five disease cohorts} against seven literature-standard methods and two state-of-the-art learners with corrected repeated cross-validation, a reproduction of the strongest published benchmark protocol, and ablations of every component (Sections~\ref{sec:setup}--\ref{sec:results}).
\item \textbf{A deployed web service} that maps the five clinical input forms to the trained ensembles and returns a risk score for each endpoint, released with the full experimental code (Sections~\ref{sec:system} and~\ref{sec:code}).
\end{enumerate}
